\scanpage{177}
% L1–4
comme quotient de
\[
\bigotimes_i E_i \oplus \bigotimes_i E_i^0
\]
et on a une suite exacte
\[
\bigotimes E_i^0 \longrightarrow \mathop{\star}\limits_{i,\lambda} E_i \unsure{\text{---}|\longrightarrow} \bigotimes \deleted{\unsure{E}}F_i\deleted{\unsure{S|s}}\longrightarrow 0
\]
% L5–8
qui correspond au fait que les systèmes $(f,f^0)$ pour lesquels $f^0=0$, \unsure{correspondent|s'identifient} aux $f$ qui \deleted{\unsure{q|se}} factorisent par les $F_i$.

% L9–13
L'application $\bigotimes E_i^0 \mathrel{\overset{\mathrm{inj}}{\unsure{\text{---}|\longrightarrow}}} \mathop{\star}\limits_{i,\lambda} E_i$ n'est sans doute pas tjrs injective, mais je dis que le choix des scindages des extensions $E_i$ définit un inverse à gauche de $\mathrm{inj}$.

% L14–19
En effet, écrivons
\[
E_i \simeq E_i^0\oplus F_i
\]
la donnée de $u_i$, déjà connue sur $E_i^0$ comme $\lambda E_i^0$, équivaut alors à la donnée de $\deleted{\unsure{t|b}}b_i:F_i \to E_i^0$
% archive sheet number: 176
